\section{Design Overview}
\label{sec:design}

SplitZip is designed around four principles:

\textbf{Lossless by default.} The primary mode guarantees bitwise exact reconstruction of every BF16 value, including rare exponents. Exactness is critical for reproducible serving and eliminates the need for quality validation per model or task.

\textbf{GPU-native throughput.} Encode and decode run entirely on GPU via Triton kernels. The codec must not become a bottleneck---it must operate faster than the network path it accelerates.

\textbf{Transfer-pipeline integration.} The codec is designed for insertion between KV cache generation and network send/receive, compatible with Mooncake's managed-buffer API. Layer-level pipelining hides codec latency behind transfer time.

\textbf{Simplicity.} The codec uses fixed-width codes (no variable-length Huffman or ANS), static per-layer codebooks (no per-batch adaptation), and byte-aligned output (no bit-level packing). This keeps the GPU kernels simple and fast.

\textbf{At a glance.} Given a BF16 KV tensor of $n$ elements (16 bits each), SplitZip:
\begin{enumerate}[nosep,leftmargin=*]
\item Splits each BF16 value into an 8-bit exponent and an 8-bit sign+mantissa.
\item Maps each exponent to a 4-bit code via a 256-entry lookup table (top-15 exponents $\to$ codes 0--14, all others $\to$ escape code 15).
\item Packs pairs of 4-bit codes into bytes, producing $n/2$ packed bytes.
\item Stores the $n$ sign+mantissa bytes unchanged.
\item Collects escaped positions and raw exponents into a tiny side stream ($\sim$0.02\% of elements).
\end{enumerate}

The compressed output is $n/2 + n + \epsilon = 1.5n + \epsilon$ bytes versus $2n$ original, yielding a ratio of $\frac{2n}{1.5n} \approx 1.333\times$. Decoding reverses the process: unpack nibbles, look up exponents via a 16-entry decode table, apply escape corrections, and recombine with sign+mantissa bytes to reconstruct BF16 values.
